\renewcommand{\BU}{\mathbf U}
\section{\emph{A priori} and beau bounds for bounded-primitive renormalization} \label{sec:conclusions}
We will begin this section by updating Theorems \ref{def vertical teaser} and \ref{thm:length-comparison} to the setting of $\psi$-quadratic-like maps.
We will then use the updated results to prove the \emph{a priori} and ``beau'' bounds for infinitely renormalizable maps of bounded-primitive type. 
This will be first expressed in the language of hyperbolic lengths of the geodesics around the small Julia sets,
and then in the form of the official modulus of quadratic-like maps that are formed as the renormalizations. 
\subsection{Estimates of extremal width and hyperbolic length for renormalizable $\psi$-quadratic-like maps}
Let us first prove the generalization of Theorem \ref{def vertical teaser} to $\psi$-quadratic-like maps:
\begin{theorem}\label{thm:pseudo-hor-ver}
Let $(i, f)\from \BU'\to \BU$ be a primitively renormalizable $\psi$-quadratic-like-map with period $p$
(and  $\KK$ is its cycle of little Julia sets).
If $\| W_\can^{\ver+\hor}(\BU\sm \KK)\|_1 > M(p)$ then
$$
   \|W^\ver_\can(\BU\sm\KK)\|_1 \geq C^{-1} \| W^\hor_\can(\BU\sm \KK)\|_1,
$$
where $M(p)$ depends only on $p$, and $C>0$ is an absolute constant.
\end{theorem}
\begin{proof}
We let $\BU^0 = \BU$, $\BU^1 = \BU'$, $f_0 = f$, and $i_0 = i$. 
As described in Section \ref{pseudo-iterate},
we can form disks $\BU^n$ and $\psi$-quadratic-like maps $(i, f)\from \BU^{n+1} \to \BU^n$.
Moreover these all have the same quadratic-like germ,
and are hence all $p$-renormalizable.
For each $n \in \N$,
we let $K^n \subset U^n$ be the central filled-Julia set of this map,
and let $\KK^n \equiv \KK_p^n \subset K^n$ be the union of small Julia sets (of the $p$-renormalization).
by a slight abuse of notation, 
we will just denote these by $K$ and $\KK \equiv \KK_p$.
We then have 
(with the slight abuse of notation described at the end of Section \ref{psi-pl}),
\begin{enumerate}
\item
a covering $f\from \BU^{n+1} \sm f^{-1} \KK \to \BU^n \sm \KK$,
\item
an inclusion $\BU^{n+1} \sm f^{-1} \KK \hookrightarrow \BU^{n+1} \sm \KK$, and
\item
an immersion $i\from\BU^{n+1} \sm \KK \to \BU^n \sm \KK$.
\end{enumerate}
The one significant change from the quadratic-like to the $\psi$-quadratic-like setting is the third map, 
which was previously an embedding, and is now an immersion. 

Now we need only go through Section \ref{sec:hor-vert} and see what needs to be changed.
In Section \ref{domination sec} we need only update the proof of Proposition \ref{restrictions} for the $\psi$-quadratic-like setting;
the proof of the updated Theorem \ref{loss of hor weight} is then identical. 
So we need only verify Properties (i) and (ii) of Proposition \ref{restrictions}.
Property (ii) follows just as before, 
because we use only the covering $f\from \BU^{n+1} \sm f^{-1} \KK \to \BU^n \sm \KK$ and the inclusion $\BU^{n+1} \sm f^{-1} \KK \hookrightarrow \BU^{n+1} \sm \KK$.
For Property (i) we need only observe that the immersion $i\from\BU^{n+1} \sm \KK \to \BU^n \sm \KK$ is still proper on the ends corresponding to $\KK$,
and Lemma \ref{e} applies to immersions as well as embeddings.

For Lemma \ref{increase of the total},
we can replace the inclusion $U^n \sm \KK \subset U \sm \KK$ 
with the immersion $i_n\from \BU^n \sm \KK \to \BU \sm \KK$;
this is still homotopic to a homeomorphism, 
so we can apply Theorem \ref{thm:total-increase} as before.
In a similar vein,
examining the proof of Theorem \ref{def vertical teaser} in Section \ref{covering-application},
we find that we need only replace $f^n\from U^n\to U$ with $f^n\from \BU^n \to \BU$,
where now this $f^n$ is the iteration described at the end of Section \ref{psi-pl}.
\end{proof}
We can now state and prove the version of Theorem \ref{thm:length-comparison} for a primitively $p$-renormalizable $\psi$-quadratic-like map $\ff \equiv (i, f)\from \BU' \to \BU$,
with $K$, $\KK$ defined as before, and geodesics $\gamma$ around $K$ and $\gamma_p$ around the central small Julia set in $\BU \sm \KK$:
\begin{theorem} \label{thm:pseudo-lengths}
Suppose that $\ff$ is a $\psi$-quadratic-like map that is primitively $p$-renormalizable.
Then, with the notation as described just above,
\begin{equation} \label{eq:pseudo-length-comparison}
|\gamma_p|> M(p) \implies |\gamma| \ge \epsilon p\, |\gamma_p|,
\end{equation}
where $M(p)$ depends only on $p$, and $\epsilon > 0$ is an absolute constant. 
\end{theorem}

\begin{proof}
This follows from Theorem \ref{thm:pseudo-hor-ver} in exactly\footnote{The proof of \eqref{eq:compare} (as Theorem 9.3 of \cite{McM}) can be readily adapted to $\psi$-quadratic-like maps.}
 the same way that Theorem \ref{thm:length-comparison} follows from Theorem \ref{def vertical teaser}. 
\end{proof}

The following is one of the main results of this paper. 
It is cited as Theorem 9.1 (as it appeared in the first arXiv version)
in \cite{decorations} and \cite{molecules}.


\begin{theorem}\label{improving of length}
For any $C \in \R^+$, there exists $\underline p$ such that for any $\bar p\geq \underline p$,
there exists $L(\bar p)$ with the following property.
Let $(i, f)\from \BU'\to \BU$ be primitively renormalizable $\psi$-quadratic-like map
with period $p$ such that  $\underline p\leq p\leq \bar p$.
Let $(j, g)\from \BV' \to \BV$ be the canonical renormalization of $f$ of period $p$.
Then (with the same notation as Theorem \ref{thm:iterated-lengths})
$$
   |\gamma(g)| > L(\bar p)\imply |\gamma(f)| > C |\gamma(g)|.
$$
\end{theorem}

We can now use Theorem \ref{thm:pseudo-lengths} to prove Theorem \ref{improving of length}.
\begin{proof}[Proof of Theorem \ref{improving of length}]
Given $C$, 
we let $\underline p = \lceil C/\epsilon \rceil$, 
where $\epsilon > 0$ is given by Theorem \ref{thm:pseudo-lengths}.
Then,
given $\bar p$,
we let $L(\bar p)$ be the maximum of $M(p)$ in Theorem \ref{thm:pseudo-lengths},
taken over all $p$ with $\underline p \le p \le \bar p$. 
Then
given the primitively $p$-renormalizable map $\ff$,
and letting $\gamma$ and $\gamma_p$ be as in Theorem \ref{eq:pseudo-length-comparison},
we have
\begin{align*} \label{eq:new-comparision}
|\gamma_p|> L(\bar p)& \implies 
|\gamma_p|> M(p) \implies 
|\gamma| \ge \epsilon p\, |\gamma_p|\\ 
& \implies
|\gamma| \ge \epsilon \underline p\, |\gamma_p| \implies
|\gamma| \ge C |\gamma_p|,
\end{align*}
where we have used \eqref{eq:pseudo-length-comparison} and our definitions of $L(\bar p)$ and $\underline p$.
Theorem \ref{improving of length} then follows from this and the first part of Theorem \ref{thm:iterated-lengths}.
\end{proof}




\subsection{\emph{A priori} and beau bounds for bounded-primitive renormalization}
We can now prove our ``if it's bad now, it was worse earlier'' statement for geodesics around the central small Julia sets.

We will have the following ``standard assumptions'' for this subsection.
We let $\ff$ be a $\psi$-quadratic-like map that is primitively renormalizable of periods $p_1\divides p_2\divides \ldots$,
with $1 < p_{k+1} / p_k \le B$ for all $k$.
The sequence of periods may be infinite,
or it may be finite,
in which case the results apply only when they refer to the $k$ for which $p_k$ exists. 
As explained in Section \ref{intro:defs},
we let $\gamma_{(k)}$ be the geodesic around $K_{p_k}(0)$ in $\C \sm \KK_{p_k}$;
we also let $\gamma_{(0)} = \gamma(\ff)$.
(If $\ff$ is a quadratic polynomial, then $\gamma_0$ does not exist, but we can let $| \gamma_0 | = 0$ and all results still apply.)

We can now prove a precise and general version of \eqref{eq:bad-now-worse-earlier}.
\begin{theorem} \label{thm:bnwe}
There exists $l$ such that for all $B$ there exists $M$:
Suppose that $\ff$ satisfies our standard assumptions.
Then,
for all $n$,
\begin{equation} \label{eq:bnwe}
|\gamma_{(n+l)}| > M \implies |\gamma_{(n)}| > 2|\gamma_{(n+l)}|.
\end{equation}
\end{theorem}

\begin{proof}
Choose $l$ such that $2^l \ge \underline p$, 
where $\underline p$ is given by Theorem \ref{improving of length} when we take $C := 2$.
Let $\bar p = B^l$,
and let $M = L(\bar p)$ in Theorem \ref{improving of length}. 
We claim that \eqref{eq:bnwe} holds for these choices (and all $n$).
To see this,
suppose that $|\gamma_{(n+l)}| > M$,
and let $\bgg \equiv (i, g)\from \BU' \to \BU$ be the canonical renormalization of $\ff$ of period $p_n$.
Let $q = p_{n+l}/p_n$;
then $\bgg$ is $q$-renormalizable,
and by Theorem \ref{thm:iterated-lengths},
we have 
\begin{equation} \label{eq:from-it}
|\gamma_{(n)}| = |\gamma(\bgg)| \text{ and } |\gamma_q(\bgg)| \ge |\gamma_{(n+l)}| > M.
\end{equation}
Since $\bar p \ge q \ge 2^l \ge \underline p$,
we can then apply Theorem \ref{improving of length} to $\bgg$ to obtain
$$
|\gamma(\bgg)| \ge 2 |\gamma_q(\bgg)|, 
$$
which combined with \eqref{eq:from-it} yields \eqref{eq:bnwe}.
\end{proof}

As a complement to Theorem \ref{thm:bnwe},
we observe that we can control the result of bounded-period renormalization. 
\begin{lemma} \label{lem:local-bound}
For all $P$ and $m > 0$ there is an $M \equiv M(m, P)$ such that 
if $\ff$ is a primitively renormalizable $\psi$-quadratic-like map of period $p \le P$ with $\gamma(\ff) \le m$,
then $|\gamma_p(\ff)| \le M$.
\end{lemma}
\begin{proof}
By Theorem \ref{thm:pseudo-lengths},
we have either $|\gamma_p| \le M(P)$ or $| \gamma_p | \le Pm /\epsilon$ for a universal $\epsilon$.
\end{proof}

Dennis Sullivan \cite{S} invented the term ``beau'', meaning \emph{bounded and eventually universally}\footnote{``beau'' looks and sounds a lot better than ``baeu''}. 
This is an extension and improvement on the \emph{a priori} bounds,
as the reader shall see in a moment. 
Here is the beau bounds for bounded-primitive renormalization, 
expressed in terms of the geodesics $\gamma_k$ around the central small Julia sets.
\begin{theorem} \label{thm:psi-beau}
For all $B$ there exists $M_0$ such that for all $m$ there exists $n_1$ and $M_1$:
Suppose $\ff$ satisfies our standard assumptions, and $| \gamma(\ff) | \le m$.
Then,
\begin{enumerate}
\item
for all $n\in \N$, $|\gamma_{(n)}| \le M_1$, and
\item
for $n \ge n_1$, $|\gamma_{(n)}| \le M_0$. 
\end{enumerate}
\end{theorem}

\newcommand{\mlong}{M_{\text{long}}}
\newcommand{\mshort}{M_{\text{short}}}

The first statement says that $| \gamma_{(n)} |$ is bounded in terms of $| \gamma(\ff) |$ (and $B$),
while the second statement says that $| \gamma_{(n)} |$ is eventually universally bounded (depending only on $B$),
with the number of renormalizations needed to become universally bounded depending on $| \gamma(\ff) |$. 

\begin{proof}
We take $l$ as in Theorem \ref{thm:bnwe},
and given $B$,
we let $\mlong$ be the $M(B)$ in Theorem \ref{thm:bnwe} 
and let $\mshort$ be $M(m, B^l)$ in Lemma \ref{lem:local-bound}.
($\mlong$ controls the long-term behavior, while $\mshort$ controls the short term.)

Let us first now Statement (1).
We take $M_1$ to be the greater value of $\mlong$ and $\mshort$.
Now suppose there were an $n$ with $|\gamma_{(n)}| > M_1$;
we take the least such value of $n$.
If $n \le l$ then $p_n \le B^l$,
and we contradict Lemma \ref{lem:local-bound}.
If $n > l$, 
then $|\gamma_{{(n-l)}}| \ge 2|\gamma_{(n)}|$ by Theorem \ref{thm:bnwe},
so $n$ was not minimal. 
This concludes the proof of (1).

For Statement (2) we take $M_0$ to be $\mlong$, which depends only on $B$.
We observe that we can rewrite Theorem \ref{thm:bnwe} as
\begin{equation} \label{eq:better}
| \gamma_{(n + l)} | \le \max(M_0, \frac12 | \gamma_{(n)} |).
\end{equation}
We let $n_1 = \ceil{1 + \log_2(M_0/\mshort)}\cdot l$. 
Now given $n \ge n_1$,
we can write $n = a + b l$,
where $a \le l$ and $b \ge \log_2(M_0/\mshort)$.
By Lemma \ref{lem:local-bound},
we have $| \gamma_{(a)} | \le \mshort$;
a simple induction with \eqref{eq:better} then implies that $| \gamma_{(a + kl)} | \le \max(M_0, 2^{-k} \mshort)$.
Letting $k = b$, we obtain our desired bound.
\end{proof}

\subsection{Quadratic polynomials and quadratic-like renormalization}
In this final subsection we return to the setting of quadratic polynomials and quadratic-like maps,
in order to derive the statements that will be more familiar (and sometimes more uto the reader. 

It will be convenient to define a \emph{semi-quadratic-like map} (with connected Julia set) to be a $\psi$-quadratic like map where the immersion is an embedding. In this case we can write it as $f\from U \to V$,
where $f$ is proper of degree 2 and $U \subset V$,
and the filled-Julia set $K \subset U$ is an FJ-set with $f^{-1}(K) = K$. 
Both quadratic-like maps and quadratic polynomials are semi-quadratic-like,
and indeed the main application of this terminology is to include both of these cases. 

In the case where $f$ is a quadratic polynomial,
the following is proven as Theorem 10.3 in \cite{McM} in the case where $M$ is sufficiently small,
and as Theorem 4.10 in \cite{towers} in general (for $f$ a with a slightly different statement that readily implies ours).
The proof here is fundamentally the same, but written with the machinery of the canonical renormalization.
\begin{theorem} \label{thm:gamma-to-modulus}
Suppose $f\from V' \to V$ is a primitively $p$-renomalizable semi-quadratic-like map,
and $|\gamma_p| \le M$.
Then $f$ has a (Douady-Hubbard) renormalization of period $p$ and modulus at least $m > 0$, for $m \equiv m(M)$. 
\end{theorem}
\begin{proof}
Let $V_p \subset V$ be the (open) disk bounded by $\gamma_p$;
and let $V'_p$ be $V_p$ pulled back along $\KK_p$,
so $V'_p$ is the unique component of $f^{-p}V_p$ that contains $K_p$.
Then $f^p\from V'_p \to V_p$ is a degree 2 cover;
we claim that $\cl{V'_p} \subset V_p$,
and that $\mod(V_p, V'_p) > m(M)$;
this would then complete the proof of the Theorem.
We let $A_p = V_p \sm K_p$ and $A'_p = V'_p \sm K_p$;
we need only show the outer boundary $\dout A'_p$  of $A'_p$ lies in $A_p$,
with a definite modulus between $\dout A'_p$ and $\dout A_p$.  

Now let $(i, g)\from \BU_p' \to \BU_p$ be the canonical renormalization of $f$ of period $p$.
We will also denote by $A_p$ and $A'_p$ their lifts to $\BU_p$ and $\BU'_p$ respectively. 
Then $g\from A'_p \to A_p$ is a degree 2 covering,
and $i\from A'_p \to A_p$ is an immersion that is degree 1 on the inner ideal boundary of $\di A'_p$
(see Remark \ref{ai:extends} and the proof of Theorem \ref{ql extension}). 
The inclusion $i(A'_p) \subset A_p$ (and lower bound on separating modulus) then follows from \eqref{ai:s} in Lemma \ref{lem:ann-immerse}.
\end{proof}
 
We can now state and prove the beau bounds for (semi-)quadratic-like maps and quadratic-like renormalization.
\begin{theorem} \label{thm:a-priori-bounds}
For all $B$ there exists $M_0$ such that for all $m$ there exists $n_1$ and $M_1$:
Suppose that $f\from U' \to U$ is a semi-quadratic-like-map that satisfies our standard assumptions, and $\abs{\mod(V \sm K)} \ge m$.
Then,
\begin{enumerate}
\item
for all $n\in \N$, 
we can find a quadratic-like $p_n$-renormalization $f^{p_n}\from V' \to V$ with $\mod(V\sm V') \ge M_1$, and 
\item
for $n \ge n_1$, 
we can find a quadratic-like $p_n$-renormalization $f^{p_n}\from V' \to V$ with $\mod(V\sm V') \ge M_0$.
\end{enumerate}
\end{theorem}
\begin{proof}
This follows immediately from Theorems \ref{thm:psi-beau} and \ref{thm:gamma-to-modulus}.
\end{proof}
Applied to quadratic-like maps, this is the standard statement of the beau bounds.

We conclude with a complete and explicit statement of the \emph{a priori} bounds,
which is an immediate corollary of Statement (1) of Theorem \ref{thm:a-priori-bounds}.
\begin{theorem}
For all $B$ there exists $m > 0$:
Let $f\from \C\to \C$ be an infinitely $B$-bounded primitively renormalizable quadratic polynomial,
with primitive periods $p_1 \divides p_2 \ldots$. 
Then for all $n$, there is a quadratic-like $p_n$-renormalization $f^{p_n}\from V' \to V$ with $\mod(V \sm V') \ge m$. 
\end{theorem}
 
